% Modelo para artigos da Sociedade Brasileira de Computação (SBC).
% Fonte do modelo: https://www.overleaf.com/latex/templates/sbc-conferences-template/blbxwjwzdngr
\documentclass[12pt]{article}

\usepackage{sbc-template}
\usepackage{graphicx,url}
\usepackage[T1]{fontenc}
\usepackage[utf8]{inputenc}
\usepackage[english,brazil]{babel}
\usepackage[hidelinks]{hyperref}

\sloppy

\title{Raça, etnia e renda regional na previsão de\\
aprovação de financiamento imobiliário}

\author{João Gabriel de Faria Beserra\inst{1}, Lucas Alves de Melo Lins\inst{1}}

\address{Instituto Federal de Educação, Ciência e Tecnologia de São Paulo (IFSP)\\
  Campus Caraguatatuba -- Avenida Bahia, 1739 -- Caraguatatuba, SP, Brasil
  % \email{Inserir os e-mails institucionais confirmados pelos autores.}
}

\begin{document}

\maketitle

\begin{resumo}
Este estudo propõe comparar regressão logística e floresta aleatória para
prever decisões de aprovação de financiamento imobiliário nos Estados Unidos,
com dados públicos do Home Mortgage Disclosure Act. Serão testados quatro
conjuntos de variáveis para avaliar como a raça, a etnia e a renda da área
onde fica o imóvel influenciam as previsões dos modelos. A análise comparará
o desempenho geral e as taxas de erro por grupo demográfico e faixa de renda local.
\end{resumo}

\section{Introdução}

Os dados públicos do \textit{Home Mortgage Disclosure Act} (HMDA) reúnem
registros de solicitações e operações de financiamento imobiliário
informados por instituições financeiras nos Estados Unidos. Esses
registros incluem a renda do solicitante, seu comprometimento com dívidas,
o valor do financiamento, características do imóvel e a decisão da
instituição~\cite{ffiec:lar}.

A partir desses dados, o estudo investiga se, além do perfil financeiro
do solicitante e das características do financiamento, a raça e a etnia
registradas e a renda da região onde fica o imóvel ajudam a explicar
quais pedidos foram aprovados ou recusados. Para isso, serão comparados
modelos que incluem informações de raça e etnia, informações de renda
regional ou ambas, tendo como referência um modelo com os dados
financeiros e as características do financiamento e do imóvel.

A pergunta central que queremos responder é como a inclusão de raça e etnia registradas e da renda
relativa do setor censitário do imóvel, separadamente e em conjunto, alteram a
previsão ou aprovação da distribuição dos erros entre grupos. Para responder
a essa pergunta, serão comparadas quatro configurações de atributos com 4 modelos diferentes e a mesma amostra de teste. A aplicação é a análise de decisões
históricas de crédito.

\section{Metodologia}

\subsection{Dados e definição de aprovação}

A unidade de análise será a solicitação de financiamento. A edição pública
de 2024 do HMDA será a referência para o recorte inicial, voltado à compra
de residência principal~\cite{ffiec:2024}. O campo \texttt{action\_taken}
define o desfecho. Os códigos 1 (financiamento contratado) e 2 (pedido
aprovado, mas não aceito) formam a classe de aprovação. O código 3
identifica as recusas. Os demais desfechos ficam fora da
análise~\cite{ffiec:lar}.

\subsection{Variáveis e treinamento}

A configuração de referência ($A$) inclui renda, idade e endividamento do
solicitante, dados do financiamento e do imóvel, instituição e ano.
A configuração $A+B$ acrescenta raça e etnia registradas no HMDA.
A configuração $A+C$ acrescenta a renda relativa da área onde fica o
imóvel, medida pela razão entre a renda familiar mediana local e a da
região de referência. A configuração $A+B+C$ reúne os dois grupos de
variáveis~\cite{ffiec:lar,ffiec:censo}. Juros, custos finais e motivos de
recusa serão excluídos por poderem revelar o desfecho.

Cada configuração será testada com regressão logística e floresta
aleatória. A divisão estratificada reservará cerca de 80\% dos registros
para treinamento e 20\% para teste, com a mesma partição em todas as
comparações. O tratamento de valores ausentes, a codificação, o ajuste
dos modelos e a escolha do limiar serão feitos na etapa de treinamento
e validação.

\subsection{Plano de avaliação}

O desempenho geral será comparado por ROC-AUC, PR-AUC para recusas e
acurácia balanceada. Por grupo demográfico e faixa de renda local,
serão informados o número de pedidos, as taxas de aprovação observada
e prevista e as taxas de falsos positivos e falsos negativos. A
comparação principal aplicará, para cada modelo, o mesmo limiar a
todos os grupos.

\section{Resultados e discussão}

\emph{A preencher após o experimento com o recorte final, o número de
solicitações e a proporção de aprovações e recusas.}

\begin{table}[htbp]
\centering
\caption{Desempenho das configurações no conjunto de teste}
\label{tab:resultados}
\small
\begin{tabular}{@{}llccc@{}}
\hline
\textbf{Algoritmo} & \textbf{Config.} & \textbf{ROC-AUC} &
\textbf{PR-AUC (recusa)} & \textbf{Acur. bal.} \\
\hline
Regressão logística & $A$       & -- & -- & -- \\
Regressão logística & $A+B$     & -- & -- & -- \\
Regressão logística & $A+C$     & -- & -- & -- \\
Regressão logística & $A+B+C$   & -- & -- & -- \\
Floresta aleatória  & $A$       & -- & -- & -- \\
Floresta aleatória  & $A+B$     & -- & -- & -- \\
Floresta aleatória  & $A+C$     & -- & -- & -- \\
Floresta aleatória  & $A+B+C$   & -- & -- & -- \\
\hline
\end{tabular}
\end{table}

% Compare cada configuração com A. Destaque as diferenças nas taxas
% de erro entre grupos e informe quantos pedidos há nos grupos citados.

A base pública oferece informações parciais sobre o histórico de crédito
e os critérios internos das instituições. A renda local se refere à
região do imóvel. As diferenças observadas serão interpretadas como
associações no recorte estudado. As estimativas de grupos pequenos
exigem cautela.

1Este estudo propõe avaliar se a inclusão de raça e etnia registradas e
da renda relativa do setor censitário do imóvel altera o desempenho
dos modelos na previsão de decisões históricas de financiamento
imobiliário. A comparação das quatro configurações na mesma amostra
de teste permitirá medir a contribuição preditiva desses atributos,
separadamente e em conjunto, em relação à configuração de referência.

A análise das taxas de falsos positivos e falsos negativos por grupo
demográfico e faixa de renda local permitirá verificar se diferenças
no desempenho geral são acompanhadas por mudanças na distribuição dos
erros. Como os dados públicos do HMDA cobrem apenas parte do histórico
de crédito e dos critérios adotados pelas instituições, a interpretação
dos resultados se limitará a associações no recorte estudado. A
execução dos experimentos permitirá estimar a magnitude dessas
diferenças e avaliar em quais grupos elas se concentram.
\begingroup
% O estilo SBC já imprime autor e ano no início de cada referência.
\makeatletter
\renewcommand{\@biblabel}[1]{}
\makeatother
\bibliographystyle{sbc}
\bibliography{referencias}
\endgroup

\end{document}
